\documentclass[10pt,a4paper]{amsart}
\usepackage[margin=19mm]{geometry}
\usepackage{amsmath,amssymb,mathtools}
\usepackage[scr=rsfs,cal=euler]{mathalfa}
\usepackage{xfrac}
\usepackage[unicode,hidelinks,bookmarks=false]{hyperref}
\newcommand{\ie}{i.e.\ }
\newcommand{\cf}{cf.\ }
\newcommand{\resp}{respectively\ }
\newcommand{\Subsection}{\subsection}
\providecommand{\nobreakdash}{\nobreak\hbox{-}}
\setlength{\emergencystretch}{2em}
\multlinegap0pt
\usepackage{amssymb}
\usepackage{mathtools}
\usepackage{xcolor}
\usepackage{tikz}
\newtheorem{proposition}{Proposition}
\newtheorem{lemma}{Lemma}
\DeclareMathOperator{\Aut}{Aut}
\title{On totally synchronizing graphs}
\author{Daniele D'Angeli$^{\ast}$, Emanuele Rodaro$^{\S}$}
\date{Source: arXiv:2607.17335v1 (CC BY 4.0)}
\begin{document}
\maketitle

\noindent\emph{Source and licence.} On totally synchronizing graphs. Version arXiv:2607.17335v1 (CC BY 4.0), distributed by arXiv under the Creative Commons Attribution 4.0 International licence, http://creativecommons.org/licenses/by/4.0/, as recorded in the arXiv licence metadata for this exact version and shown on the abstract page https://arxiv.org/abs/2607.17335v1. Full licence notice: \texttt{licenses/arxiv-2607.17335-CC-BY-4.0.txt}.\par\smallskip

\noindent\textit{Editorial scope.} Counted unit: the article's lemma lem:triangular-walks-lift (source0.tex lines 939--952). The article's complete original proof of it is delivered with it (source0.tex lines 954--1028, one \texttt{proof} environment of 75 lines). No mathematics is written, repaired or re-derived: the statement and the proof are the article's own lines, in the article's own notation, and no retained line is substituted or re-flowed.  Retained uncounted as the source-local context the proof needs: lines 804--808.  Nothing was omitted from any retained span, so the package carries no omission record.  No retained line contains a citation, so the wrapper carries no bibliography and no bibliography entry.  No retained line contains a figure environment, an image-inclusion command or drawing code, so no image is delivered and the package declares no asset dependency.  Editorial formatting only, in addition: an AMS \texttt{amsart} preamble that restates the article's own declarations and macros used by the retained lines, namely the environments \texttt{proposition}, \texttt{lemma} on separate counters, together with the standard packages those lines need.  The retained environments print in this document as Proposition 1, Lemma 1 in order of appearance, whereas the article prints the same items with the numbering of its own full document.  The complete native source of the article as distributed in the arXiv e-print for this exact version is bundled unchanged as \texttt{sources/205560/source0.tex}.
\begin{proposition}\label{prop:vertex-rigid-refinement-permutation}
Let $G=(V,E,s,t)$ be a finite strongly connected $k$-out graph, with $k\ge 2$.
Assume that there exists a permutation $\tau:V\to V$ such that $|\delta_u(\tau(u))|\ge 1$ for every $u\in V$, and such that the graph $G^-=(V,E^-,s,t)$, obtained by removing exactly one edge from $\delta_u(\tau(u))$ for each $u\in V$, is still strongly connected ($(k-1)$-out graph). Then there exist a finite strongly connected $k$-out graph $\widetilde G=(\widetilde V,\widetilde E,s,t)$ and a congruence $\rho_0$ on $\widetilde G$ such that $\widetilde G/\rho_0 \simeq G$, 
and every automorphism of $\widetilde G$ fixes every vertex. Equivalently, $\Aut(\widetilde G)=\Aut_E(\widetilde G)$.
\end{proposition}

\begin{lemma}\label{lem:triangular-walks-lift}
Under the assumptions of
Proposition~\ref{prop:vertex-rigid-refinement-permutation},
assume moreover that there exist a vertex $v_0\in V$ and closed walks
$D_1,\dots,D_r$ in $G$, all based at $v_0$, such that, for each $i$,
the edge $f_i$ occurs exactly once in $D_i$ and does not occur in any
of the walks $D_1,\dots,D_{i-1}$. Then the integers $L_u$ in the construction of
Proposition~\ref{prop:vertex-rigid-refinement-permutation} may be chosen
so that the resulting refinement $\widetilde G$ contains closed walks
$\widetilde D_1,\dots,\widetilde D_r$, all based at $r_{v_0}$, such
that, for each $i$, the edge $\widetilde f_i$ occurs exactly once in
$\widetilde D_i$ and does not occur in any of the walks
$\widetilde D_1,\dots,\widetilde D_{i-1}$.
\end{lemma}

\begin{proof}
We choose the refinement in
Proposition~\ref{prop:vertex-rigid-refinement-permutation}
as follows. For each $u\in V$, let $o(u)$ be the length of the
$\tau$-cycle containing $u$, and choose pairwise distinct positive
integers $L_u$ such that
\[
o(u)\mid L_u
\qquad (u\in V).
\]
Such a choice is possible by taking sufficiently large distinct
multiples of the corresponding cycle lengths. Then
\[
\tau^{L_u}(u)=u,
\qquad
\tau^{L_u+1}(u)=\tau(u).
\]
Hence the coding walk associated with $u$ is a walk from $r_u$ to
$r_{\tau(u)}$:
\[
r_u\longrightarrow c_{u,1}\longrightarrow\cdots
\longrightarrow c_{u,L_u}\longrightarrow r_{\tau(u)}.
\]
Let $\widetilde G$ be the refinement obtained from this choice. For each $u,w\in V$, assign to every edge $e\in\delta_u(w)$ a walk
$\widehat e$ in $\widetilde G$ from $r_u$ to $r_w$ as follows. If $w\neq\tau(u)$, there are exactly
\[
m(u,w)=|\delta_u(w)|
\]
core edges from $r_u$ to $r_w$; choose a bijection between
$\delta_u(w)$ and these core edges.

If $w=\tau(u)$, there are
\[
m(u,\tau(u))-1
\]
core edges from $r_u$ to $r_{\tau(u)}$, together with the coding walk
\[
r_u\longrightarrow c_{u,1}\longrightarrow\cdots
\longrightarrow c_{u,L_u}\longrightarrow r_{\tau(u)}.
\]
Since
\[
m(u,\tau(u))
=
\bigl(m(u,\tau(u))-1\bigr)+1,
\]
choose a bijection between $\delta_u(\tau(u))$ and the set consisting
of these core edges and this coding walk.
Thus every edge $e\in E(G)$ is assigned a unique walk $\widehat e$ in
$\widetilde G$ from $r_{s(e)}$ to $r_{t(e)}$. For each $i$, let
$\widetilde D_i$ be obtained from $D_i$ by replacing every edge $e$
with the corresponding walk $\widehat e$. Since $D_i$ is closed and
based at $v_0$, the walk $\widetilde D_i$ is closed and based at
$r_{v_0}$.
Fix $i$. By hypothesis, the edge $f_i$ occurs exactly once in $D_i$
and does not occur in $D_1,\dots,D_{i-1}$. If $\widehat f_i$ is a core edge, set
\[
\widetilde f_i:=\widehat f_i.
\]
Since $f_i$ occurs exactly once in $D_i$, the edge
$\widetilde f_i$ occurs exactly once in $\widetilde D_i$.
By injectivity of the assignment $e\mapsto\widehat e$, it does not
occur in any of $\widetilde D_1,\dots,\widetilde D_{i-1}$.
If $\widehat f_i$ is the coding walk associated with some $u\in V$, set
\[
\widetilde f_i:=
\bigl(r_u\longrightarrow c_{u,1}\bigr).
\]
This edge belongs to no assigned walk other than $\widehat f_i$.
Since $f_i$ occurs exactly once in $D_i$, the edge
$\widetilde f_i$ occurs exactly once in $\widetilde D_i$.
Since $f_i$ does not occur in any of $D_1,\dots,D_{i-1}$,
the edge $\widetilde f_i$ does not occur in any of
$\widetilde D_1,\dots,\widetilde D_{i-1}$. Hence, for each $i$, the edge $\widetilde f_i$ occurs exactly once in $\widetilde D_i$ and does not occur in any of $\widetilde D_1,\dots,\widetilde D_{i-1}$.
\end{proof}

\end{document}
